\chapter{Субстрат: типизированный темпоральный граф}
\label{ch:substrate}

\section{Что такое узел и что такое ребро}

Память Рин --- направленный граф $G = (V, E)$, в котором узлы несут содержание
(воспоминания, факты, сущности, концепты), а рёбра --- типизированные и
знаковые отношения между ними. Ключевая мысль главы: и узлы, и рёбра несут
метаданные, и именно метаданные, а не сам факт связи, делают возможными все
дальнейшие механизмы.

\begin{definition}[Узел памяти]
Узел $v \in V$ --- это запись с содержанием $x_v$ и набором полей состояния,
над которыми работают операторы памяти. Узел одновременно является и
``данными'' (что помнится), и ``убеждением'' (насколько это считается
истинным), и ``следом'' (насколько активно и свежо).
\end{definition}

\subsection{Схема узла}

\begin{longtable}{p{0.24\textwidth} p{0.10\textwidth} p{0.56\textwidth}}
\toprule
\textbf{Поле} & \textbf{Тип} & \textbf{Смысл и кто с ним работает} \\
\midrule
\endhead
\texttt{id} & int & идентификатор \\
\texttt{content} & text & содержание в естественном виде \\
\texttt{embedding} $x_v$ & vec$[d]$ & векторное представление; retrieval-релевантность, мост баз \\
\texttt{node\_kind} & enum & вид узла (observation\slash claim\slash hypothesis\slash preference\slash \dots); задаёт смысл \texttt{confidence} (см.\ раздел~\ref{sec:nodekind}) \\
\texttt{confidence} $c_v$ & $[-1, 1]$ & для truth-apt узлов --- знаковое убеждение ($+$да/$-$нет/$0$хз); для остальных --- надёжность оценки (раздел~\ref{sec:nodekind}) \\
\texttt{valence} & vec$[2]$ & валентность--возбуждение; сцепка с настроением \\
\texttt{layer}\slash\texttt{ts} & enum\slash time & масштаб (сессия\dots жизнь) и/или метка времени; консолидация \\
\texttt{provenance} & list & источники и/или обоснования; TMS, ревизия \\
\texttt{need\_to\_*} & flags & ``исследовать\slash пересмотреть''; любопытство, скраб \\
\texttt{importance} & $[0, 1]$ & значимость: ставится salience при кодировании (гл.~\ref{ch:personality}), подкрепляется центральностью; гейтит прунинг (гл.~\ref{ch:timescales}) \\
\texttt{domains} & tags & членство в доменах знаний; несёт importance по членству, а не диффузией (гл.~\ref{ch:timescales}) \\
\texttt{activation} $a_v$ & float & транзиентная активация; обход \\
\texttt{degree}\slash\texttt{centrality} & float & производные структурные; топология \\
\bottomrule
\end{longtable}

\subsection{Схема ребра}

Ребро $(u \to v)$ несёт направленность (\texttt{src}, \texttt{dst}), тип
отношения \texttt{type} (\textsc{isA}, \textsc{partOf}, \textsc{supports},
\textsc{contradicts}, \textsc{associates}, \dots) и два независимых числовых
канала: вес перехода $w_{uv} \ge 0$ (нормируется по узлу, для обхода) и знак
$s_{uv} \in \{+1, -1\}$ (полярность отношения, для пропагации уверенности).
Знак --- хранимый лейбл, внутреннее свойство связи: ``эта связь за или
против?'', а не вычисляемая величина. Его нельзя ``посчитать по пути'', он
известен при создании ребра. Изменение уверенности $\Delta \conf$ --- уже
производное, оно вычисляется в пропагации с использованием знака
(глава~\ref{ch:propagation}).

\begin{intuition}[Веса --- про свидетельство, а не про характер]
Забегая вперёд (глава~\ref{ch:personality}): вес ребра выражает
\emph{эвиденциальную} силу связи и уверенность, а не ``характер''. Соблазн
задать личность статическими множителями на рёбра-реакции
(``грусть$\to$сарказм $\times 2{,}0$'') мы разберём и отвергнем: характер
должен входить как смещение во время обхода от текущего состояния, а не как
вшитый в граф множитель.
\end{intuition}

\section{Почему рёбра направленные, типизированные и знаковые}

Плоская модель ``узлы и связи одинакового сорта'' ломается сразу на трёх
вещах: направленность (асимметрия причинности и времени), тип (нельзя валить
covalent-связь, близость, ``подтверждает'' и ``является частью'' в один
линейный канал) и знак (без него не отличить ``всё усилилось'' от ``всё
перевернулось'').

Правильный оператор поэтому не общий лапласиан сглаживания, а
\emph{relation-specific} передача сообщения. Сообщение по ребру:
\begin{equation}
m_{u \to v} = \phi_{r_{uv}}\!\big(h_u,\, h_v,\, e_{uv},\, \Delta t,\, c_{uv},\, s_{uv}\big),
\label{eq:message}
\end{equation}
где $r_{uv}$ --- тип отношения, $\Delta t$ --- возраст или временная
дистанция, $c_{uv}$ --- уверенность, $s_{uv}$ --- источник или эпистемический
статус. Затем узел агрегирует входящие сообщения --- и здесь важно не одну
статистику, а набор (сумма, среднее, максимум, дисперсия по типам отношений),
потому что для памяти существенны и кратность подтверждений, и пропорции, и
наличие экстремального свидетельства.

\begin{workedbox}[Почему одна статистика теряет смысл]
Ключевое свойство \emph{среднего} --- инвариантность к дублированию. Пусть про
факт есть свидетельства $\{+, -\}$ (один за, один против) и в другом случае
$\{+, +, -, -\}$ (двое за, двое против). Считая $+ = 1$, $- = -1$, получаем
\[
\mathrm{mean}\{+, -\} = 0 = \mathrm{mean}\{+, +, -, -\}.
\]
Среднее одинаково, хотя во втором случае свидетелей вдвое больше. Для памяти
это принципиально: ``четверо разделились'' --- не то же, что ``один на один''.
\emph{Сумма} же сохраняет кратность, а \emph{максимум} по знаку --- наличие
сильного контрсвидетельства. Отсюда правило: агрегатор хранит \emph{набор}
статистик (сумма, среднее, максимум, дисперсия), а не одну, --- иначе
теряется либо кратность, либо пропорция, либо экстремум.
\end{workedbox}

\begin{intuition}[Аналогия с электрической цепью]
Полезно думать о графе как о цепи с проводниками разного сорта.
Направленность --- это диоды; тип ребра --- материал проводника со своей
проводимостью; знак --- это полярность. Плоская модель --- цепь, где все
элементы заменены одинаковыми резисторами: она проведёт ток, но потеряет всю
логику схемы.
\end{intuition}

\begin{center}
\begin{tikzpicture}[>=Stealth,font=\footnotesize,
  mem/.style={draw, rounded corners, minimum height=7mm, inner sep=3pt},
  inf/.style={draw, rounded corners, fill=blue!5, minimum height=7mm, inner sep=3pt}]
  \node[mem] (a) {разбил кружку};
  \node[mem, below=9mm of a] (d) {юзер хотел новую};
  \node[inf, right=26mm of a] (r) {реакция: грусть};
  \node[inf, right=22mm of d] (x) {вывод: уместно облегчение};
  \draw[->,thick] (a) -- node[above,font=\scriptsize]{$+$ xReact} (r);
  \draw[->,dashed] (a) -- (x);
  \draw[->,dashed] (d) -- (x);
  \draw[->,thick,red!60!black] (x) -- node[right,font=\scriptsize]{$-$ ослабляет} (r);
\end{tikzpicture}
\captionof{figure}{Противоречие не самостоятельное, а \emph{контекстное}. Базовое $+$xReact ``разбил $\to$ грусть'' ослабляется не узлом ``хотел новую'' самим по себе, а совместным выводом ``разбил $+$ хотел новую $\to$ уместно облегчение'' (пунктир --- провенанс; это $n$-арная конфигурация, глава~\ref{ch:fusion}). Именно она делает возможным уместное ``поздравляю''.}
\label{fig:cup}
\end{center}

\section{Виды узлов и что значит уверенность}
\label{sec:nodekind}

Раньше мы говорили: ``узел --- это одновременно данные и убеждение''. Красиво,
но для реализации грубо: не всякий узел --- утверждение об истинности.
``Португалия имеет такой-то режим магазинов'', ``этот вывод поддержан
источником'', ``гипотеза House / Cast Away ложна'' --- да, это truth-apt
утверждения, и для них $\conf \in [-1, 1]$ означает истинность. Но
``Ярославу заходит такой подкол'', ``Рин не любит такой тон'', ``этот человек
устал'', ``у Анны есть кот'' --- это предпочтения, реакции, склонности,
состояния. ``Уверенность 0.8, что Ярослав любит $X$'' --- не то же самое, что
``вероятность истинности факта 0.9''. Предпочтения меняются, зависят от
контекста, имеют исключения.

Поэтому у узла есть поле \texttt{node\_kind}, и оно задаёт смысл
\texttt{confidence}:

\begin{center}
\begin{tabular}{p{0.22\textwidth} p{0.28\textwidth} p{0.28\textwidth}}
\toprule
\textbf{node\_kind} & \textbf{что это} & \textbf{смысл confidence} \\
\midrule
observation & прямое наблюдение & истинность (высокая) \\
claim & утверждение о мире & истинность \\
hypothesis & вывод или гипотеза & истинность (низкая, провизорно) \\
episode & конкретное событие & было / не было \\
preference & вкус или склонность & надёжность оценки, не истинность \\
reaction\_pattern & ``на $X$ реагирую $Y$'' & надёжность, Beta-постериор \\
relationship\_state & состояние отношения / кто-что-знает & надёжность модели \\
boundary & граница или нежелание & устойчивость \\
skill/procedure & навык / как-делать & освоенность \\
self\_model & убеждение о себе & надёжность самооценки \\
mood\_summary & свёртка настроения & это метка, не убеждение \\
\bottomrule
\end{tabular}
\end{center}

Для не-truth-apt видов confidence читается как надёжность или устойчивость
оценки, а не ``вероятность истинности''. Предпочтения и реакции лучше хранить
как модель с условиями, а не как голый скаляр истинности.

\begin{principle}[Confidence priors = source\_trust $\times$ kind, а не плоская константа]
Импортированным узлам, то есть эмбриональному графу (ConceptNet, ATOMIC),
нельзя давать плоскую уверенность на всё. Baseline должен зависеть от двух
осей:
\[
\mathrm{confidence}(\mathrm{node}) =
\mathrm{source\_trust}(\mathrm{source}) \times
\mathrm{baseline}(\mathrm{node\_kind}).
\]
По \texttt{kind}: ConceptNet-концепт --- truth-apt claim, стартует выше;
ATOMIC-эпизод --- социально-каузальная оценка реакции, стартует ниже, потому
что это надёжность мягкого прогноза, а не истинность. По источнику: разные
базы заслуживают разного доверия.
\end{principle}

Два принципиальных решения здесь особенно важны.

\begin{itemize}
\item Источник --- это провенанс, а не метка. \texttt{source} уже несётся в
провенансе (\texttt{/c/en/...} vs \texttt{atomic}); не надо плодить
эпистемические метки \texttt{prior\_conceptnet} / \texttt{prior\_atomic}. Метка
\texttt{prior} --- это эпистемический класс для packer'а, а источник ---
ортогональные метаданные для калибровки.
\item Не выводить node-confidence из веса ребра. Вес ConceptNet --- это
корроборация отношения, канал ребра, а не истинность узла. Вес гонит
диффузию, confidence узла --- про доверие к содержанию.
\end{itemize}

\begin{intuition}[Почему знаковый скаляр, а не тройка мнения]
Напрашивается более богатое представление --- тройка субъективной логики
\emph{belief / disbelief / uncertainty}, где неопределённость несёт собственную
координату. Мы сознательно берём один знаковый скаляр $c \in [-1, 1]$
(эквивалент $2P - 1$: $+1$ ``да'', $-1$ ``нет'', $0$ ``не знаю''). Причины
две. Во-первых, скаляр уже разводит неверие ($c < 0$) и незнание
($c \approx 0$). Во-вторых, третью ось --- ``сколько за этим свидетельств /
насколько оценка устойчива'' --- мы и так несём отдельно и разнотипно: через
\texttt{node\_kind}, через \texttt{need\_to\_discover}, через
Beta-неопределённость реакций (гл.~\ref{ch:personality}) и провенанс.
Неопределённость у нас не потеряна, она просто не запихнута в один кортеж.
\end{intuition}

\begin{intuition}[Кто-что-знает --- тоже вид убеждения]
Особо отметим \texttt{relationship\_state} в смысле моделирования чужих
знаний. ``Рин знает $X$'', ``Ярослав знает, что Анна сказала $X$'', ``Анна не в
курсе, что Рин это запомнила'' --- это разные узлы: убеждения Рин о знаниях
других. Это не про приватность-как-политику, а про то, что theory-of-mind ---
законный и важный вид памяти: без него Рин не отличит ``я это знаю'' от ``он
думает, что я этого не знаю''.
\end{intuition}

\subsection{От словаря к поведению: операционная матрица}
\label{sec:kindmatrix}

Список видов --- это словарь; реализация спросит \emph{поведение}. Каждый
\texttt{node\_kind} обязан ответить на семь вопросов: какие операторы его
читают, какие могут в него писать, как считается confidence, как работает
распад, какие типы рёбер допустимы, можно ли его fusion-ить и можно ли
показывать LLM в context-packer'е (гл.~\ref{ch:runtime}).

Полную матрицу $11 \times 7$ заполнять сейчас --- значит выдумывать точность,
которой нет. Поэтому ниже --- каркас с наиболее решающими столбцами и
уверенными клетками; прочее помечено как TBD.

\begin{center}
\begingroup
\small
\setlength{\tabcolsep}{4pt}
\begin{tabular}{@{}p{0.18\textwidth} p{0.21\textwidth} p{0.16\textwidth} p{0.13\textwidth} p{0.16\textwidth}@{}}
\toprule
\textbf{kind} & \textbf{смысл confidence} & \textbf{распад} & \textbf{fusion?} & \textbf{в packer?} \\
\midrule
observation & истинность & медленный & да (посылка) & да \\
claim & истинность & медленный & да & да \\
hypothesis & истинность (низкая) & быстрый, если не подтверждена & да & да, с меткой ``гипотеза'' \\
preference & надёжность оценки & по использованию & осторож\-но & да \\
\shortstack[l]{reaction\_\\pattern} & Beta-надёжность & по успеху / провалу & нет, это модель & как контекст \\
\shortstack[l]{relationship\_\\state} & надёжность модели & медленный & осторож\-но & \shortstack[l]{по ситуации\\(TBD)} \\
boundary & устойчивость & очень медленный & нет & да \\
\shortstack[l]{self\_\\model} & надёжность самооценки & медленный & да (метакогниция) & да \\
\shortstack[l]{mood\_\\summary} & это метка & транзиент\-но & нет & редко \\
\bottomrule
\end{tabular}
\endgroup
\end{center}

Столбцы ``кто читает / кто пишет / допустимые рёбра'' определяются там, где
живёт соответствующий оператор, а не дублируются здесь. Матрица --- это
\emph{контракт} поведения на тип, и она сознательно неполна: заполняется по
мере постройки.

\section{Провенанс и обоснование}
\label{sec:prov}

Поле \texttt{provenance} шире, чем ``источник''. Для \emph{исходных} фактов
провенанс --- откуда факт узнан. Для узлов, \emph{порождённых выводом}
(fusion), провенанс должен хранить \emph{обоснование} --- множество
родительских узлов, из которых вывод получен. Эта деталь оказывается
критической в главе~\ref{ch:epistemics}: именно наличие обоснований позволяет,
опровергнув посылку, автоматически пометить на пересмотр её следствия. Без
записанных обоснований опровержение посылки не находит потомков --- и
рождается ``ползучая вера''.

\section{Что значит «коммит» и кто его гейтит}

Важнейшая граница схемы --- между активным рабочим убеждением и записью. Не
всё, что всплыло или сгенерировалось, попадает в граф как факт.

\begin{principle}[Гипотеза, а не факт]
Всё, что порождено выводом или генерацией (fusion, реакция, ассоциация),
пишется провизорно: с низкой уверенностью, с провенансом-обоснованием и, при
необходимости, с флагом \texttt{need\_to\_discover}. Повышение до ``записи'',
то есть переход в более медленный слой, происходит только через накопление
подтверждающих свидетельств. Управление --- на границе коммита.
\end{principle}

\begin{lstlisting}[style=pseudo]
def commit(G, item, kind):
    if kind == "observed":          # прямое наблюдение
        c0 = source_trust(item.source)
    if kind == "inferred":          # порождено fusion / выводом
        c0 = LOW                    # гипотеза, не факт
        item.provenance = item.parents      # обоснование для TMS
    node = G.add_node(item, confidence=c0, layer="session")
    propagate(G, node)              # см. главу про лавину
    return node
\end{lstlisting}

\section{Кража схемы у Wikidata}

Есть готовая, обкатанная модель данных, уже решившая ровно те задачи, что
встают перед метаданными Рин, --- модель утверждений \emph{Wikidata}. Брать
её \emph{содержание} для существа, мимикрирующего под человека, не нужно:
человек помнит мало фактов, его сила в ассоциативности. Но её \emph{схему}
стоит перенять целиком:

\begin{itemize}
\item квалификаторы (точка во времени, начало, конец) --- темпоральная
валидность: факт истинен в интервале, а не вечно;
\item references --- провенанс и обоснование;
\item ранг (\emph{preferred} / \emph{normal} / \emph{deprecated}) ---
недеструктивная ревизия: устаревшее не удаляют, а депрецируют, помечая новое
предпочтительным; история сохраняется для машинерии поддержки истинности.
\end{itemize}

Иными словами, факт стоит моделировать как
\[
\text{(утверждение)} + \text{(интервал валидности)} +
\text{(провенанс)} + \text{(статус / ранг)},
\]
и это буквально схема Wikidata. Не изобретать метаданные заново, а взять
проверенный шаблон.

\section{Врождённое и приобретённое: не граница, а континуум}

Соблазнительно разделить узлы на ``врождённые константы'' (эмбриональный
граф) и ``приобретённые'' (нажитые, с временными метками). Как первое
приближение полезно, но как жёсткая граница неверно. Консолидация
\emph{двигает} записи через границу (эпизод $\to$ обобщённое знание), а
ревизия должна уметь \emph{понижать} и ``врождённое'' --- иначе унаследованную
ошибку здравого смысла Рин не исправит никогда. Правильная картина ---
\emph{континуум глубины консолидации}, где порог опрокидывания убеждения
растёт с глубиной, но нигде не становится бесконечным. К оси
``пластичность -- упрямство'' мы ещё не раз вернёмся
(главы~\ref{ch:propagation}, \ref{ch:epistemics}).
